\subsection{Concentration for monotone sub-additive functions}

The main step in our proof of Theorem~\ref{thm:random_tensor} is a type of concentration inequality for monotone sub-additive functions on the hypercube.

In order to obtain such a result we will consider a notion of \emph{two-point distance} from sets $A\subseteq \bset{n}$, which intuitively measures how many coordinates of some given point $x$ cannot be captured by any two elements of~$A$.

\begin{definition}
Given a point $x\in \bset{n}$ and a set $A \subseteq \bset{n}$, the two-point distance between $x$ and $A$ is
\beqn
h_2(x; A) = \min_{y, z \in A} \big|\big\{i \in [n] \st x_i \neq y_i \And x_i \neq z_i\big\}\big|.
\eeqn
\end{definition}

Note that this ``distance'' $h_2(x; A)$ can be zero even if $x \notin A$;
for instance, $h_2\big(00;\, \{01, 10\}\big) = 0$.
The reason for considering this notion is that one gets much better concentration for $h_2$ than one does for the usual Hamming distance.
This is shown by the following result, which is a special case of an inequality of Talagrand~\cite[Theorem 3.1.1]{Talagrand:1995}:

\begin{theorem}\label{thm:Talagrand}
For any parameter $\sigma\in [0,1]$ and set $A \subseteq \bset{n}$, we have that
\beqn
\Pr_{x\sim \pi_\sigma^n}\big[h_2(x; A) \geq k\big] \leq 2^{-k} \pi_\sigma^n(A)^{-2} \quad \text{for all $k\geq 0$.}
\eeqn
\end{theorem}

The next lemma is the key result needed for proving our main theorem for tensors, and it deals more abstractly with monotone, sub-additive Lipschitz functions on the hypercube.
We endow~$\bset{n}$ with the usual partial order, where $x\leq y$ if the support of~$x$ is contained in the support of~$y$.
A function $f:\bset{n} \to \R$ is monotone if $f(x) \leq f(y)$ whenever $x\leq y$, and it is sub-additive if $f(x+y) \leq f(x) + f(y)$ for all $x,y\in\bset{n}$ with  disjoint supports.
Finally,~$f$ is $1$-Lipschitz if
\beqn
\big|f(x) - f(y)\big| \leq \big|\big\{i \in [n] \st x_i \neq y_i\big\}\big| \quad \text{for all $x, y \in \bset{n}$}.
\eeqn
For a string $x\in \bset{n}$ and set $I\subseteq[n]$, we let~$x_I\in \bset{n}$ be the string that equals~$x$ on the indices in~$I$ and is zero elsewhere.
The lemma that follows and its proof are inspired by an argument of Schechtman~\cite[Corollary~12]{Schechtman2003}.

\begin{lemma}[Concentration inequality] \label{lem:subadditive}
Let $f:\bset{n} \to \R_+$ be a monotone, sub-additive $1$-Lipschitz function such that $f({\bf 0}) =0$ and $f({\bf 1}) = r$.
Then,
\begin{align*}
    \Pr_{x\sim \pi_{\sigma}^n} \big[f(x) \leq \sigma r/4\big] &\leq \sqrt{\frac{3}{2\sigma}}\, 2^{-\sigma r/12} \quad &\text{if $\,0 < \sigma < 1/2$,} \\
    \Pr_{x\sim \pi_{\sigma}^n} \big[f(x) \leq r/6\big] &\leq \sqrt{2}\, 2^{-r/12} &\text{if $\,1/2 \leq \sigma \leq 1$.}
\end{align*}
\end{lemma}

\begin{proof}
We first prove the first inequality above.
Let $\sigma\in (0, 1/2)$, and consider the set
\beqn
A = \big\{ x\in \bset{n} \st f(x)\leq \sigma r/4 \big\}.
\eeqn
By Talagrand's Inequality (Theorem~\ref{thm:Talagrand}) we have
\beqn
\Pr_{x\sim \pi_{\sigma}^n}\big[ h_2(x; A)\geq \sigma r/6\big] \leq 2^{-\sigma r/6} \,\Pr_{x\sim \pi_{\sigma}^n}\big[ f(x)\leq \sigma r/4\big]^{-2},
\eeqn
so to finish the proof it suffices to show that
\beqn
\Pr_{x\sim \pi_{\sigma}^n}\big[ h_2(x; A)\geq \sigma r/6\big] \geq 2\sigma/3.
\eeqn

We claim that
\beq \label{eq:h2inclusion}
\big\{ x\in \bset{n} \st f(x)\geq 2\sigma r/3 \big\} \subseteq \big\{ x\in \bset{n} \st h_2(x; A)\geq \sigma r/6 \big\}.
\eeq
Indeed, if $h_2(x; A) < \sigma r/6$ then there are $y, z\in A$ such that
\beqn
\big|\big\{i \in [n] \st x_i \neq y_i \And x_i \neq z_i\big\}\big| < \sigma r/6.
\eeqn
Denote $I = \big\{i \in [n] \st x_i \neq y_i \And x_i \neq z_i\big\}$, $J = \big\{i \in [n] \st x_i = y_i\big\}$ and $J' = \big\{i \in [n] \st x_i \neq y_i \And x_i = z_i\big\}$;
note that these sets partition~$[n]$, and by assumption $|I| < \sigma r/6$.
We then have
\begin{align*}
    f(x) &\leq f(x_I + x_J) + f(x_{J'}) \\
    &= f(x_I + y_J) + f(z_{J'}) \\
    &\leq |I| + f(y_J) + f(z_{J'}) \\
    &\leq |I| + f(y) + f(z) \\
    &< 2\sigma r/3,
\end{align*}
so $\big\{h_2(x; A) < \sigma r/6\big\} \subseteq \big\{f(x) < 2\sigma r/3\big\}$ and inclusion \eqref{eq:h2inclusion} follows.
It thus suffices to show that $\Pr_{x\sim \pi_{\sigma}^n}\big[ f(x)\geq 2\sigma r/3\big] \geq 2\sigma/3$.

We will next prove that, for any integer $k\geq 1$, we have
\beq \label{eq:coupling}
\Pr_{x\sim \pi_{1/k}^n}\big[ f(x)\geq r/k\big] \geq 1/k.
\eeq
This is done by a simple \emph{coupling argument}, which will be important for us again later on.
Consider a uniformly random ordered $k$-partition $(I_1, \dots, I_k)$ of $[n]$;
thus the $k$ sets $I_i$ are pairwise disjoint and have union $[n]$, with each of the $k^n$ possible such $k$-tuples having the same probability.
Denoting by $\1_I$ the indicator function of set $I$, we have
\beqn
r = f(\1_{[n]}) = f(\1_{I_1} + \dots + \1_{I_k}) \leq \sum_{i=1}^k f(\1_{I_i}),
\eeqn
so $f(\1_{I_i}) \geq r/k$ holds for at least one $i\in [k]$ in \emph{every} ordered partition $(I_1, \dots, I_k)$.
By symmetry, it follows that $\Pr \big[f(\1_{I_1}) \geq r/k \big] \geq 1/k$.
Since the marginal distribution of $\1_{I_1}$ (and every other $\1_{I_i}$) is precisely $\pi_{1/k}^n$, we conclude that
\begin{equation*}
\Pr_{x\sim \pi_{1/k}^n} \big[f(x) \geq r/k\big] = \Pr \big[f(\1_{I_1}) \geq r/k \big] \geq 1/k,
\end{equation*}
as wished.

Now let $k = \lceil 1/\sigma \rceil$.
Since $0 < \sigma < 1/2$, we have $2\sigma/3 \leq 1/k \leq \sigma$, and so by monotonicity of $f$
\begin{align*}
    \Pr_{x\sim \pi_{\sigma}^n}\big[ f(x)\geq 2\sigma r/3\big] &\geq \Pr_{x\sim \pi_{1/k}^n}\big[ f(x)\geq 2\sigma r/3\big] \\
    &\geq \Pr_{x\sim \pi_{1/k}^n}\big[ f(x)\geq r/k\big].
\end{align*}
From inequality~\eqref{eq:coupling} we then conclude that
\beqn
\Pr_{x\sim \pi_{\sigma}^n}\big[ f(x)\geq 2\sigma r/3\big] \geq 1/k \geq 2\sigma/3,
\eeqn
finishing the proof of the first inequality in the statement of the lemma.

The second inequality is proven using the same arguments, now applied to the parameter $\sigma = 1/2$ and the set $A = \big\{ x\in \bset{n} \st f(x)\leq r/6 \big\}$.
By the same argument as above,
\begin{equation*}
    \big\{x\in \bset{n}\st f(x) \geq r/2\big\} \subseteq \big\{x\in \bset{n}\st h_2(x;A) \geq r/6\big\}.
\end{equation*}
Talagrand's inequality and coupling imply that
\begin{align*}
    2^{-r/6}\Pr_{x\sim\pi_{1/2}^n}\big[f(x) \leq r/6\big]^{-2} &\geq \Pr_{x\sim \pi_{1/2}^n}\big[h_2(x;A) \geq r/6\big]\\
    &\geq \Pr_{x\sim\pi_{1/2}^n}\big[f(x) \geq r/2\big]\\
    &\geq \frac{1}{2}.
\end{align*}
Rearranging then gives the result for~$\sigma =  1/2$;
we obtain slightly better bounds since in this case $1/\sigma = 2$ is an integer, and so no rounding errors occur.
By monotonicity of $f$, the same bound continues to hold for all $\sigma > 1/2$.
\end{proof}


\subsection{The Random Restriction Theorem}

It is now a simple matter to use Lemma~\ref{lem:subadditive} to prove the Random Restriction Theorem for tensors.

\begin{proof}[ of Theorem~\ref{thm:random_tensor}]
Let $\rank: (\F^{\infty})^{\otimes d} \to \R$ be a natural rank function, and $T \in \bigotimes_{i=1}^d \F^{[n_i]}$ be a tensor.
We wish to show that
\beqn
\Pr_{I_1\sim [n_1]_{\sigma}, \dots, I_d\sim [n_d]_{\sigma}}\big[ \rank \big(T_{|I_{[d]}}\big) \geq \kappa(\sigma) \cdot \rank(T)\big] \geq 1 - C(\sigma) e^{-\kappa(\sigma) \rank(T)}
\eeqn
for some well-chosen constants $C(\sigma), \kappa(\sigma) > 0$ depending only on the value of $\sigma > 0$ and the order $d$ of the tensor.
We consider here the case where $\sigma < 1/2$, as the case where $\sigma \geq 1/2$ is analogous.\footnote{This second case is also an immediate consequence of first case together with monotonicity of $\rank$ under restrictions.}

Define the function $f_1: \bset{n_1} \to \R_+$ as follows:
for $x \in \bset{n_1}$ with support $J$, $f_1(x)$ is the rank of the subtensor of $T$ whose indices of the first row belong to $J$.
In formula:
\beqn
f_1(\1_{J}) = \rank\big( T_{|J \times \prod_{j=2}^d [n_j]} \big) \quad \text{for all } J\subseteq [n_1].
\eeqn
By the definition of natural rank, $f_1$ is a monotone, sub-additive $1$-Lipschitz function with maximum value $\rank(T)$.
Since $\1_J \sim \pi_{\sigma}^{n_1}$ when $J\sim [n_1]_{\sigma}$, it follows from Lemma~\ref{lem:subadditive} that
\beq \label{eq:f1}
\Pr_{I_1\sim [n_1]_{\sigma}}\bigg[ \rank\big( T_{|I_1 \times \prod_{j=2}^d [n_j]} \big) \geq \frac{\sigma \rank(T)}{4}\bigg] \geq 1 - \sqrt{\frac{3}{2\sigma}}\, 2^{-\frac{\sigma}{12} \rank(T)}.
\eeq

For any fixed $I_1 \subseteq [n_1]$ satisfying $f_1(\1_{I_1}) \geq \sigma \rank(T)/4$, we define the function $f_2: \bset{n_2} \to \R_+$ by
\beqn
f_2(\1_{J}) = \rank\big( T_{|I_1 \times J \times \prod_{j=3}^d [n_j]} \big) \quad \text{for all } J\subseteq [n_2].
\eeqn
This function is again monotone, sub-additive and $1$-Lipschitz, and it has maximum value at least~$\sigma \rank(T)/4$.
By Lemma~\ref{lem:subadditive} we have
\beqn
\Pr_{I_2\sim [n_2]_{\sigma}}\bigg[ \rank\big( T_{|I_1 \times I_2 \times \prod_{j=3}^d [n_j]} \big) \geq \Big(\frac{\sigma}{4}\Big)^2 \rank(T)\bigg] \geq 1 - \sqrt{\frac{3}{2\sigma}}\, 2^{-\frac{\sigma}{12} \frac{\sigma}{4} \rank(T)}.
\eeqn
Since this holds whenever $f_1(\1_{I_1}) \geq \sigma \rank(T)/4$, taking the union bound together with inequality~\eqref{eq:f1} we obtain
\beqn
\Pr_{I_1\sim [n_1]_{\sigma}, I_2\sim [n_2]_{\sigma}}\bigg[ \rank\big( T_{|I_1 \times I_2 \times \prod_{j=3}^d [n_j]} \big) \geq \Big(\frac{\sigma}{4}\Big)^2 \rank(T)\bigg] \geq 1 - 2\sqrt{\frac{3}{2\sigma}}\, 2^{-\frac{\sigma}{12} \frac{\sigma}{4} \rank(T)}.
\eeqn
Proceeding in this same way for each row of the tensor, and always taking the union bound, we eventually conclude that
\beqn
\Pr_{I_1\sim [n_1]_{\sigma}, \dots, I_d\sim [n_d]_{\sigma}}\bigg[ \rank \big(T_{|I_{[d]}}\big) \geq \Big(\frac{\sigma}{4}\Big)^d \rank(T)\bigg] \geq 1 - d\sqrt{\frac{3}{2\sigma}}\, 2^{-\frac{\sigma}{12} (\frac{\sigma}{4})^{d-1} \rank(T)},
\eeqn
which finishes the proof.
\end{proof}

